% ============================================================
% Chapitre 2 : R\'egression Lin\'eaire et Polynomiale
% ============================================================
\chapter{R\'egression Lin\'eaire et Polynomiale}
\label{ch:regression}

\begin{intuition}[title=\textbf{Id\'ee centrale}]
La r\'egression lin\'eaire est le mod\`ele fondamental de l'apprentissage supervis\'e : on cherche
la <<~meilleure~>> relation affine entre des variables explicatives $\bm{x}$ et une variable
cible $y \in \R$, au sens des moindres carr\'es.
\end{intuition}

% ============================================================
\section{R\'egression lin\'eaire simple}
\label{sec:reg-simple}

\subsection{Mod\`ele}

On dispose de $n$ observations $(x_i, y_i)_{i=1}^{n}$ avec $x_i \in \R$ et $y_i \in \R$.
On suppose le mod\`ele :
\[
  y_i = \beta_0 + \beta_1 x_i + \varepsilon_i, \qquad \varepsilon_i \sim \mathcal{N}(0, \sigma^2),
  \quad \varepsilon_i \text{ i.i.d.}
\]

\begin{definition}[Fonction de co\^ut des moindres carr\'es]
La fonction de co\^ut \emph{Ordinary Least Squares} (OLS) est :
\[
  \mathcal{L}(\beta_0, \beta_1) = \sum_{i=1}^{n} (y_i - \beta_0 - \beta_1 x_i)^2
\]
L'objectif est de trouver $(\hat{\beta}_0, \hat{\beta}_1) = \argmin_{\beta_0, \beta_1} \mathcal{L}(\beta_0, \beta_1)$.
\end{definition}

\subsection{D\'erivation des estimateurs OLS}

\begin{theoreme}[Estimateurs des moindres carr\'es]
\label{thm:ols-simple}
Les estimateurs OLS de la r\'egression lin\'eaire simple sont :
\[
  \hat{\beta}_1 = \frac{\sum_{i=1}^{n}(x_i - \bar{x})(y_i - \bar{y})}{\sum_{i=1}^{n}(x_i - \bar{x})^2}
  = \frac{\Cov(X, Y)}{\Var(X)}, \qquad
  \hat{\beta}_0 = \bar{y} - \hat{\beta}_1 \bar{x}
\]
o\`u $\bar{x} = \frac{1}{n}\sum_{i} x_i$ et $\bar{y} = \frac{1}{n}\sum_{i} y_i$.
\end{theoreme}

\begin{proof}
On annule les d\'eriv\'ees partielles de $\mathcal{L}$ :
\[
  \frac{\partial \mathcal{L}}{\partial \beta_0} = -2\sum_{i=1}^{n}(y_i - \beta_0 - \beta_1 x_i) = 0
  \implies \hat{\beta}_0 = \bar{y} - \hat{\beta}_1 \bar{x}
\]
\[
  \frac{\partial \mathcal{L}}{\partial \beta_1} = -2\sum_{i=1}^{n} x_i(y_i - \beta_0 - \beta_1 x_i) = 0
\]
En substituant $\hat{\beta}_0$ dans la seconde \'equation :
\[
  \sum_{i=1}^{n} x_i \bigl(y_i - \bar{y} + \hat{\beta}_1(\bar{x} - x_i)\bigr) = 0
  \implies \hat{\beta}_1 = \frac{\sum_{i}(x_i - \bar{x})(y_i - \bar{y})}{\sum_{i}(x_i - \bar{x})^2}
\]
\end{proof}

\begin{formules}[title=\textbf{Coefficient de d\'etermination}]
\[
  R^2 = 1 - \frac{\sum_{i}(y_i - \hat{y}_i)^2}{\sum_{i}(y_i - \bar{y})^2}
  = 1 - \frac{\text{RSS}}{\text{TSS}} \in [0, 1]
\]
$R^2$ mesure la proportion de variance expliqu\'ee par le mod\`ele.
\end{formules}

% ============================================================
\section{R\'egression lin\'eaire multiple}
\label{sec:reg-multiple}

\subsection{Formulation matricielle}

On g\'en\'eralise \`a $p$ variables explicatives. Pour $n$ observations :
\[
  \bm{y} = \bm{X}\bm{\beta} + \bm{\varepsilon}
\]
o\`u $\bm{X} \in \R^{n \times (p+1)}$ est la matrice de design (premi\`ere colonne de 1),
$\bm{\beta} \in \R^{p+1}$ et $\bm{\varepsilon} \sim \mathcal{N}(\bm{0}, \sigma^2 \bm{I}_n)$.

\begin{definition}[Fonction de co\^ut matricielle]
\[
  \mathcal{L}(\bm{\beta}) = \norm{\bm{y} - \bm{X}\bm{\beta}}^2
  = (\bm{y} - \bm{X}\bm{\beta})^\top (\bm{y} - \bm{X}\bm{\beta})
\]
\end{definition}

\subsection{\'Equations normales}

\begin{theoreme}[\'Equations normales]
\label{thm:normal-equations}
Si $\bm{X}^\top \bm{X}$ est inversible, l'unique minimiseur de $\mathcal{L}$ est :
\[
  \hat{\bm{\beta}} = (\bm{X}^\top \bm{X})^{-1} \bm{X}^\top \bm{y}
\]
\end{theoreme}

\begin{proof}
On d\'eveloppe et on d\'erive :
\begin{align*}
  \mathcal{L}(\bm{\beta})
  &= \bm{y}^\top \bm{y} - 2\bm{\beta}^\top \bm{X}^\top \bm{y}
     + \bm{\beta}^\top \bm{X}^\top \bm{X} \bm{\beta} \\
  \nabla_{\bm{\beta}} \mathcal{L}
  &= -2\bm{X}^\top \bm{y} + 2\bm{X}^\top \bm{X}\bm{\beta} = \bm{0}
\end{align*}
D'o\`u $\bm{X}^\top \bm{X} \hat{\bm{\beta}} = \bm{X}^\top \bm{y}$. La matrice hessienne
$2\bm{X}^\top \bm{X}$ est semi-d\'efinie positive, donc c'est bien un minimum.
\end{proof}

\subsection{Interpr\'etation g\'eom\'etrique}

\begin{intuition}[title=\textbf{Projection orthogonale}]
Le vecteur $\hat{\bm{y}} = \bm{X}\hat{\bm{\beta}}$ est la \textbf{projection orthogonale}
de $\bm{y}$ sur l'espace colonne de $\bm{X}$, not\'e $\text{Col}(\bm{X})$.
La matrice de projection est $\bm{H} = \bm{X}(\bm{X}^\top \bm{X})^{-1}\bm{X}^\top$
(<<~hat matrix~>>). Le vecteur des r\'esidus $\hat{\bm{\varepsilon}} = \bm{y} - \hat{\bm{y}}$
est orthogonal \`a $\text{Col}(\bm{X})$.
\end{intuition}

\begin{proposition}[Propri\'et\'es de la hat matrix]
\label{prop:hat-matrix}
La matrice $\bm{H} = \bm{X}(\bm{X}^\top \bm{X})^{-1}\bm{X}^\top$ v\'erifie :
\begin{enumerate}
  \item $\bm{H}^2 = \bm{H}$ (idempotente),
  \item $\bm{H}^\top = \bm{H}$ (sym\'etrique),
  \item $\text{tr}(\bm{H}) = p + 1$.
\end{enumerate}
\end{proposition}

% --- TikZ : droite de r\'egression et r\'esidus ---
\begin{figure}[htbp]
\centering
\begin{tikzpicture}[scale=1.1]
  % axes
  \draw[->] (0,0) -- (7,0) node[right]{$x$};
  \draw[->] (0,0) -- (0,5) node[above]{$y$};
  % droite de r\'egression
  \draw[thick, blue!70!black] (0.3,0.8) -- (6.5,4.5)
    node[right, font=\small]{$\hat{y} = \hat{\beta}_0 + \hat{\beta}_1 x$};
  % points et r\'esidus
  \foreach \x/\y/\yhat in {1/1.8/1.2, 2/1.5/1.8, 2.5/2.8/2.1,
    3/2.2/2.4, 3.8/3.5/2.9, 4.5/2.8/3.3, 5/4.2/3.7, 5.8/3.8/4.0,
    6.2/4.8/4.3}{
    \fill[red!80!black] (\x, \y) circle (2pt);
    \draw[red!50, dashed, thick] (\x, \y) -- (\x, \yhat);
  }
  % l\'egende r\'esidu
  \draw[<->, red!60!black] (5.3, 3.7) -- (5.3, 4.2);
  \node[right, red!60!black, font=\footnotesize] at (5.35, 3.95) {$\varepsilon_i$};
\end{tikzpicture}
\caption{Droite de r\'egression et r\'esidus (segments rouges).}
\label{fig:regression-residus}
\end{figure}

% ============================================================
\section{R\'egression polynomiale}
\label{sec:reg-poly}

\begin{definition}[Mod\`ele polynomial de degr\'e $d$]
On pose $\phi(x) = (1, x, x^2, \ldots, x^d)^\top \in \R^{d+1}$ et :
\[
  y = \bm{\phi}(x)^\top \bm{\beta} + \varepsilon
  = \beta_0 + \beta_1 x + \beta_2 x^2 + \cdots + \beta_d x^d + \varepsilon
\]
C'est un mod\`ele \textbf{lin\'eaire en les param\`etres} $\bm{\beta}$, m\^eme s'il est
non-lin\'eaire en $x$.
\end{definition}

\begin{attention}[title=\textbf{Sur-apprentissage polynomial}]
Un polyn\^ome de degr\'e $d = n - 1$ passe exactement par les $n$ points
($R^2 = 1$ sur l'ensemble d'entra\^inement), mais g\'en\'eralise tr\`es mal.
C'est un cas classique de \textbf{sur-apprentissage} (overfitting). Le choix de $d$
est un probl\`eme de s\'election de mod\`ele (cf.\ chapitre~\ref{ch:regularisation}).
\end{attention}

\begin{remarque}
En pratique, on construit la matrice de design de Vandermonde :
\[
  \bm{\Phi} = \begin{pmatrix}
    1 & x_1 & x_1^2 & \cdots & x_1^d \\
    1 & x_2 & x_2^2 & \cdots & x_2^d \\
    \vdots & & & \ddots & \vdots \\
    1 & x_n & x_n^2 & \cdots & x_n^d
  \end{pmatrix} \in \R^{n \times (d+1)}
\]
et on applique les \'equations normales : $\hat{\bm{\beta}} = (\bm{\Phi}^\top \bm{\Phi})^{-1}\bm{\Phi}^\top \bm{y}$.
\end{remarque}

% ============================================================
\section{Descente de gradient pour la r\'egression lin\'eaire}
\label{sec:gradient-descent}

\begin{algorithme}[title=\textbf{Descente de gradient (batch)}]
\textbf{Entr\'ee :} donn\'ees $(\bm{X}, \bm{y})$, taux d'apprentissage $\eta > 0$,
nombre d'it\'erations $T$.

\textbf{Initialisation :} $\bm{\beta}^{(0)} = \bm{0}$.

\textbf{Pour} $t = 0, 1, \ldots, T-1$ :
\[
  \bm{\beta}^{(t+1)} = \bm{\beta}^{(t)} - \eta \nabla_{\bm{\beta}} \mathcal{L}(\bm{\beta}^{(t)})
  = \bm{\beta}^{(t)} + \frac{2\eta}{n} \bm{X}^\top (\bm{y} - \bm{X}\bm{\beta}^{(t)})
\]

\textbf{Sortie :} $\hat{\bm{\beta}} = \bm{\beta}^{(T)}$.
\end{algorithme}

\begin{proposition}[Convergence]
Si $0 < \eta < \frac{2}{\lambda_{\max}(\bm{X}^\top \bm{X} / n)}$, la descente de gradient
converge vers la solution des \'equations normales $\hat{\bm{\beta}}_{\text{OLS}}$.
\end{proposition}

\begin{bonnepratique}[title=\textbf{Normalisation des variables}]
Avant d'appliquer la descente de gradient, il est essentiel de \textbf{centrer-r\'eduire}
les variables ($z_j = \frac{x_j - \mu_j}{\sigma_j}$) pour que toutes les composantes du
gradient aient des \'echelles comparables. Cela acc\'el\`ere la convergence.
\end{bonnepratique}

% ============================================================
\section{Impl\'ementation Python}
\label{sec:reg-python}

\subsection{Impl\'ementation depuis z\'ero}

\begin{codeblock}[title=\textbf{R\'egression lin\'eaire -- impl\'ementation manuelle}]
\begin{minted}{python}
import numpy as np

class LinearRegressionScratch:
    """Regression lineaire par equations normales."""

    def fit(self, X, y):
        # Ajout de la colonne de biais
        n = X.shape[0]
        X_b = np.c_[np.ones((n, 1)), X]
        # Equations normales
        self.beta_ = np.linalg.solve(
            X_b.T @ X_b, X_b.T @ y
        )
        return self

    def predict(self, X):
        n = X.shape[0]
        X_b = np.c_[np.ones((n, 1)), X]
        return X_b @ self.beta_

    def r_squared(self, X, y):
        y_pred = self.predict(X)
        ss_res = np.sum((y - y_pred) ** 2)
        ss_tot = np.sum((y - np.mean(y)) ** 2)
        return 1 - ss_res / ss_tot
\end{minted}
\end{codeblock}

\subsection{Descente de gradient depuis z\'ero}

\begin{codeblock}[title=\textbf{Descente de gradient pour la r\'egression lin\'eaire}]
\begin{minted}{python}
class LinearRegressionGD:
    """Regression lineaire par descente de gradient."""

    def __init__(self, lr=0.01, n_iter=1000):
        self.lr = lr
        self.n_iter = n_iter

    def fit(self, X, y):
        n, p = X.shape
        X_b = np.c_[np.ones((n, 1)), X]
        self.beta_ = np.zeros(p + 1)
        self.losses_ = []

        for _ in range(self.n_iter):
            residuals = y - X_b @ self.beta_
            gradient = -(2 / n) * X_b.T @ residuals
            self.beta_ -= self.lr * gradient
            self.losses_.append(np.mean(residuals ** 2))

        return self

    def predict(self, X):
        X_b = np.c_[np.ones((X.shape[0], 1)), X]
        return X_b @ self.beta_
\end{minted}
\end{codeblock}

\subsection{Avec scikit-learn}

\begin{codeblock}[title=\textbf{R\'egression lin\'eaire et polynomiale avec scikit-learn}]
\begin{minted}{python}
from sklearn.linear_model import LinearRegression
from sklearn.preprocessing import PolynomialFeatures
from sklearn.pipeline import make_pipeline
from sklearn.model_selection import cross_val_score
from sklearn.metrics import mean_squared_error, r2_score
import numpy as np

# --- Donnees synthetiques ---
np.random.seed(42)
X = 2 * np.random.rand(100, 1)
y = 4 + 3 * X[:, 0] + np.random.randn(100)

# --- Regression lineaire simple ---
model = LinearRegression()
model.fit(X, y)
print(f"beta_0 = {model.intercept_:.3f}")
print(f"beta_1 = {model.coef_[0]:.3f}")
print(f"R^2    = {model.score(X, y):.4f}")

# --- Regression polynomiale (degre 3) ---
poly_model = make_pipeline(
    PolynomialFeatures(degree=3, include_bias=False),
    LinearRegression()
)
poly_model.fit(X, y)

# --- Validation croisee ---
scores = cross_val_score(
    poly_model, X, y, cv=5,
    scoring='neg_mean_squared_error'
)
print(f"MSE (CV 5-folds) = {-scores.mean():.4f}")
\end{minted}
\end{codeblock}

\begin{codeoutput}
\begin{minted}{text}
beta_0 = 4.114
beta_1 = 2.840
R^2    = 0.8765
MSE (CV 5-folds) = 1.0432
\end{minted}
\end{codeoutput}

% ============================================================
\section{Propri\'et\'es statistiques des estimateurs OLS}
\label{sec:proprietes-ols}

\begin{theoreme}[Gauss--Markov]
Sous les hypoth\`eses du mod\`ele lin\'eaire ($\E[\bm{\varepsilon}] = \bm{0}$,
$\Cov(\bm{\varepsilon}) = \sigma^2 \bm{I}_n$), l'estimateur OLS
$\hat{\bm{\beta}}$ est le \textbf{meilleur estimateur lin\'eaire sans biais} (BLUE) :
pour tout estimateur lin\'eaire sans biais $\tilde{\bm{\beta}}$,
\[
  \Var(\bm{a}^\top \hat{\bm{\beta}}) \leq \Var(\bm{a}^\top \tilde{\bm{\beta}})
  \quad \forall\, \bm{a} \in \R^{p+1}
\]
\end{theoreme}

\begin{formules}[title=\textbf{Distribution des estimateurs}]
Sous les hypoth\`eses gaussiennes :
\[
  \hat{\bm{\beta}} \sim \mathcal{N}\bigl(\bm{\beta},\; \sigma^2 (\bm{X}^\top \bm{X})^{-1}\bigr)
\]
L'estimateur sans biais de $\sigma^2$ est :
\[
  \hat{\sigma}^2 = \frac{1}{n - p - 1}\norm{\bm{y} - \bm{X}\hat{\bm{\beta}}}^2
  = \frac{\text{RSS}}{n - p - 1}
\]
\end{formules}

\begin{remarque}
La d\'ecomposition biais-variance pour la r\'egression lin\'eaire donne, pour une
nouvelle observation $\bm{x}_0$ :
\[
  \E\bigl[(\hat{y}_0 - y_0)^2\bigr]
  = \underbrace{\Bias(\hat{y}_0)^2}_{= 0 \text{ si mod\`ele correct}}
  + \Var(\hat{y}_0) + \sigma^2
\]
\end{remarque}

% ============================================================
\section{Exercices}
\label{sec:reg-exercices}

\begin{exercice}[Niveau 1 -- Calcul direct]
On observe les donn\'ees suivantes :
\begin{center}
\begin{tabular}{c|ccccc}
$x_i$ & 1 & 2 & 3 & 4 & 5 \\
\hline
$y_i$ & 2.1 & 3.9 & 6.2 & 7.8 & 10.1
\end{tabular}
\end{center}
\begin{enumerate}
  \item Calculer $\bar{x}$, $\bar{y}$, $\sum(x_i - \bar{x})(y_i - \bar{y})$
        et $\sum(x_i - \bar{x})^2$.
  \item En d\'eduire $\hat{\beta}_1$ et $\hat{\beta}_0$.
  \item Calculer $R^2$ et interpr\'eter.
  \item Pr\'edire $y$ pour $x = 6$.
\end{enumerate}
\end{exercice}

\begin{exercice}[Niveau 2 -- R\'egression multiple et diagnostic]
On consid\`ere le jeu de donn\'ees \texttt{Boston Housing} (ou \'equivalent) avec $p = 5$
variables s\'electionn\'ees.
\begin{enumerate}
  \item \'Ecrire la matrice de design $\bm{X}$ et calculer $\hat{\bm{\beta}}$ avec les
        \'equations normales en Python (sans scikit-learn).
  \item Comparer avec \texttt{LinearRegression} de scikit-learn.
  \item Tracer les r\'esidus $\hat{\varepsilon}_i$ en fonction de $\hat{y}_i$.
        Que v\'erifie-t-on avec ce graphique ?
  \item R\'ealiser un Q-Q plot des r\'esidus. Commenter la normalit\'e.
  \item Calculer le VIF (\emph{Variance Inflation Factor}) de chaque variable.
        Identifier d'\'eventuelles colin\'earit\'es.
\end{enumerate}
\end{exercice}

\begin{exercice}[Niveau 3 -- Analyse th\'eorique et descente de gradient]
\begin{enumerate}
  \item D\'emontrer le th\'eor\`eme de Gauss--Markov.
  \item Montrer que pour la r\'egression polynomiale de degr\'e $d$ sur $n$ points,
        le nombre de param\`etres effectifs est $d + 1$ et que la matrice de Vandermonde
        est mal conditionn\'ee pour $d$ grand. En d\'eduire l'int\'er\^et de la
        r\'egularisation (lien avec le chapitre~\ref{ch:regularisation}).
  \item Impl\'ementer la descente de gradient stochastique (SGD) et la descente de
        gradient par mini-batchs. Comparer les courbes de convergence avec la descente
        de gradient batch.
  \item Montrer que la condition $\eta < 2/\lambda_{\max}$ est n\'ecessaire \`a la
        convergence. \emph{Indication :} analyser le syst\`eme dans la base propre de
        $\bm{X}^\top \bm{X}$.
\end{enumerate}
\end{exercice}
